%-------------------------------------------------------------------------------
% 8. 2024 부산대학교 총장배 창의비행체 경진대회 (드론 짐벌·랜딩기어)
%-------------------------------------------------------------------------------
\cvsection{8. 2024 부산대학교 총장배 창의비행체 경진대회}

\projsummary{고정익(무동력) 기체 추종 촬영을 위한 카메라 2축 짐벌 및 리트랙팅 랜딩기어 시스템 구축}
\projfigL[4.6cm]{p8.jpg}

\projpart{\faBullhorn\hspace{1.3mm}대회 및 임무 소개}
\begin{cvitems}
  \item {\textbf{임무}\enskip 이륙 $\to$ 고도 80m 이내 분리 $\to$ 고정익 기체 추종 촬영·실시간 영상 송출 $\to$ 정밀 착륙·영상 데이터 제출.}
  \item {\textbf{배경}\enskip 고도 80m 상공에서 고정익 기체를 추종 촬영하기 위해, 랜딩기어에 의한 시야 간섭을 제거하고 카메라 짐벌을 실시간 제어하는 통합 시스템이 필요.}
  \item {\textbf{목표}\enskip 1km 이상 장거리에서도 지연 없이 2축 카메라 짐벌과 리트랙팅 랜딩기어를 동시 제어하는 HC-12 기반 무선 통신 체계·통합 하드웨어 제어 로직 설계.}
  \item {\textbf{기간·장소}\enskip 2024.08.22 -- 2024.08.23 · 부산대학교 양산캠퍼스 드론 R\&D센터}
  \item {\textbf{기술 스택}\enskip \pill{Raspberry Pi} \pill{Arduino} \pill{HC-12 (RF)} \pill{OpenCV} \pill{PWM Control} \pill{SoftwareSerial}}
\end{cvitems}

\projfig[4.4cm]{p8mission.png}

\projpart{\faMagnifyingGlass\hspace{1.3mm}설계 과정 및 수행 내용}

\stephead{1}{아날로그 짐벌 제어 시스템}
\begin{substeps}
  \item 초기 w/a/s/d 키 단계별 제어는 비행 중 기체의 빠른 움직임을 부드럽게 추적하기엔 반응이 느리고 불연속.
  \item 조이스틱 모듈로 아날로그 제어 환경을 구축, 기울기를 실시간 각도(angle1·angle2)로 변환해 서보에 전달.
  \item 카메라 Tilt·Pan을 매끄럽게 제어해 기체 추종 성능 확보.
\end{substeps}

\stephead{2}{시야각(FOV) 최적화 + 리트랙팅 랜딩기어}
\begin{substeps}
  \item 고정형 랜딩기어는 짐벌이 하방을 향할 때 시야를 가려 촬영·추종 임무에 간섭.
  \item DC 모터+드라이버(L298N)로 이륙 직후 접고 착륙 시 전개하는 리트랙팅 시스템 설계.
  \item command 값(1 수납/2 전개/0 정지)으로 회전 방향을 결정 --- 비행 중 시야 100\% 개방, 착륙 시 기체 보호.
\end{substeps}

\stephead{3}{HC-12 RF 기반 장거리 무선 통신}
\begin{substeps}
  \item Wi-Fi 등 단거리 통신은 1km 이상에서 신호 감쇄가 심해 명령 지연·두절 위험.
  \item 433MHz HC-12 모듈로 지상 관제소-기체 간 End-to-End 무선 통신망 설계, SoftwareSerial로 디버깅 포트와 독립 채널 확보.
  \item 데이터 패킷(command,angle1,angle2.) 구조를 정의하고 마침표(.)를 종결 문자로 쓰는 파싱 로직 적용.
\end{substeps}

\achievement{예선 통과 및 \textbf{본선 진출}.}
